\clearpage\section*{Solutions EV1-A}\label{EV1-A-CORR}
\textbf{MCQ:} 1 B, 2 A, 3 D, 4 B, 5 A, 6 B, 7 B, 8 B, 9 D, 10 C, 11 A, 12 A. Each correct answer: 0.5 point; total 6.
\paragraph{Checks of the fluency questions.} Order of operations; sum of fractions; subtraction of exponents; multiplying coefficient 1.2; division by 0.8; product 1.25 times 0.8; solving an affine equation; vertical increment divided by horizontal; mean; angle divided by 360; complement to 1; substitution in brackets.
\subsection*{Exercise 1 -- Decreasing reserve}
\begin{enumerate}
\item $u_0=192$, $u_{n+1}=\dfrac{u_n}{2}$; common ratio $\dfrac{1}{2}$. \textit{Marking: 1 point. Result and basic justification expected; half a point possible for a correct but unfinished approach.}
\item $192,96,48,24,12$, sum 372. \textit{Marking: 2 points. Method or justification: 1; calculation and conclusion: 1.}
\item $u_4=12$, $u_5=6$: index 5. Initialise n=0,u=192; while u is at least 10, divide by 2 and increase n. \textit{Marking: 2 points. Method or justification: 1; calculation and conclusion: 1.}
\end{enumerate}
\subsection*{Exercise 2 -- Parameter and tangent}
\begin{enumerate}
\item $\Delta=4m^2-16=4(m^2-4)$. \textit{Marking: 1 point. Result and basic justification expected; half a point possible for a correct but unfinished approach.}
\item Two roots if $m<-2$ or $m>2$; one if $m=\pm2$; none if $-2<m<2$. \textit{Marking: 2 points. Method or justification: 1; calculation and conclusion: 1.}
\item For m=1, rate $2+h$, derivative 2 and $f(2)=4$; tangent $y=2x$. \textit{Marking: 2 points. Method or justification: 1; calculation and conclusion: 1.}
\end{enumerate}
\subsection*{Exercise 3 -- Locus via the midpoint}
\begin{enumerate}
\item $MI^2-16$, by scalar identities. \textit{Marking: 1 point. Result and basic justification expected; half a point possible for a correct but unfinished approach.}
\item $MI^2=25$: circle with centre I, radius 5. \textit{Marking: 1 point. Result and basic justification expected; half a point possible for a correct but unfinished approach.}
\item For -16, $MI^2=0$: I alone. For -20, $MI^2=-4$: empty. \textit{Marking: 2 points. Method or justification: 1; calculation and conclusion: 1.}
\end{enumerate}
\paragraph{Calibration.} Target duration: 20--25 min for the fluency questions, 80--85 min for the exercises, 10--15 min for proofreading. These durations are editorial estimates to be measured in class. Accept any correct method; do not penalise twice a numerical error that is propagated consistently. The proposed marking scheme is not a decision of the team to harmonise marking.
